\chapter*{Abstract (update)}
 
High-quality data is essential for effective decision-making, data analytics and training AI models. However, real-world datasets often contain errors, inconsistencies and missing values that can negatively impact their quality. This project aims to create an automated data cleaning pipeline that combines the strength of ontologies, knowledge graphs and Large Language Models (LLMs). The goal is to eliminate manual prompt engineering by using a structured system that dynamically generates prompts for data cleaning tasks based on a dataset's profile.
The core of the system is a generalized ontology-driven knowledge graph, which stores information about data types (e.g., integers, categories, dates), common data quality issues (such as missing values or incorrect formats), validation rules (like regex checks) and correction strategies. This ontology acts as the system’s decision-making guide.
When a dataset is uploaded, a data profiler analyzes each column's characteristics and identifies inter-column relationships. This information is used to query a knowledge graph, which contains structured prompt templates enriched with cleaning strategies, examples, and task-specific instructions. These templates are then instantiated with dataset-specific data and passed to an LLM to generate cleaned data. This approach makes data cleaning automatic, adaptive to different datasets and scalable. 
% Resulting in a flexible framework that enables ontology-aware data preprocessing across domains.



